\chapter{数学、动态规划与几何}

\section{快速幂与组合数}

\cardmeta{基础必会}{幂、组合数、取模}{快速幂 \(O(\log b)\)，组合数预处理 \(O(n)\)}{模数为质数且 \(n<\texttt{mod}\) 时使用阶乘逆元}

\begin{lstlisting}[style=cpp]
int qpow(int a, int b, int mod) {
    // res 保存已经处理的二进制位贡献；a 每轮平方
    int res = 1 % mod;
    a %= mod;
    while (b) {
        // 当前最低位为 1，说明这一位的幂需要乘进答案
        if (b & 1) res = res * a % mod;
        a = a * a % mod;
        b >>= 1;
    }
    return res;
}

struct Comb {
    // fac[i]=i!，ifac[i]=(i!)^{-1}，均按 mod 取模
    int mod;
    vector<int> fac, ifac;
    void init(int n, int p) {
        mod = p;
        fac.assign(n + 1, 1);
        ifac.assign(n + 1, 1);
        for (int i = 1; i <= n; ++i) fac[i] = fac[i - 1] * i % mod;
        ifac[n] = qpow(fac[n], mod - 2, mod);
        // 费马小定理求 n! 的逆元，再递推所有较小逆阶乘
        for (int i = n; i; --i) ifac[i - 1] = ifac[i] * i % mod;
    }
    int C(int n, int k) const {
        if (k < 0 || k > n) return 0;
        return fac[n] * ifac[k] % mod * ifac[n - k] % mod;
    }
};
\end{lstlisting}

\begin{warnbox}
模数不是质数时不能直接使用费马逆元；组合数下标超过模数时先查 Lucas 或非质数模数模板。
\end{warnbox}

\section{扩展欧几里得与扩展 CRT}

\cardmeta{现场常用}{多个同余式、模数不互质、求最小非负解}{每次合并 \(O(\log m)\)}{无解时必须返回标志}

\begin{lstlisting}[style=cpp]
int exgcd(int a, int b, int &x, int &y) {
    // 返回 gcd(a,b)，并求出 a*x+b*y=gcd(a,b)
    if (!b) { x = 1; y = 0; return a; }
    int x0, y0;
    int g = exgcd(b, a % b, x0, y0);
    x = y0;
    y = x0 - (a / b) * y0;
    return g;
}

// 合并 x = a1 (mod m1), x = a2 (mod m2)
bool merge_crt(int a1, int m1, int a2, int m2, int &a, int &m) {
    int x, y, g = exgcd(m1, m2, x, y);
    int d = a2 - a1;
    // 两式可合并的充要条件：差值能被 gcd(m1,m2) 整除
    if (d % g) return false;
    int lcm = m1 / g * m2;
    int k = (i128)d / g * x % (m2 / g);
    if (k < 0) k += m2 / g;
    a = (a1 + (i128)m1 * k) % lcm;
    if (a < 0) a += lcm;
    m = lcm;
    return true;
}
\end{lstlisting}

\begin{card}
\textbf{样例}\\
\(x\equiv1\pmod4,\ x\equiv3\pmod6\) 的最小非负解为 \(9\pmod{12}\)。\\
\(x\equiv0\pmod2,\ x\equiv1\pmod4\) 无解，必须输出题目规定的无解标志。
\end{card}

\section{矩阵快速幂}

\cardmeta{现场常用}{线性递推、固定维度状态转移}{\(O(k^3\log n)\)}{矩阵维度必须很小}

\begin{lstlisting}[style=cpp]
struct Mat {
    // a[i][j] 表示从状态 j 对第 i 个新状态的线性贡献
    int n;
    vector<vector<int>> a;
    Mat(int n = 0, bool identity = false): n(n), a(n, vector<int>(n)) {
        if (identity) for (int i = 0; i < n; ++i) a[i][i] = 1;
    }
};
Mat operator*(const Mat &A, const Mat &B) {
    Mat C(A.n);
    // 普通矩阵乘法；i128 防止中间乘法溢出 long long
    for (int i = 0; i < A.n; ++i)
        for (int k = 0; k < A.n; ++k) if (A.a[i][k])
            for (int j = 0; j < A.n; ++j)
                C.a[i][j] = (C.a[i][j] +
                    (i128)A.a[i][k] * B.a[k][j]) % 1000000007;
    return C;
}
Mat mpow(Mat A, int e) {
    // R 初始为单位矩阵，按 e 的二进制位做快速幂
    Mat R(A.n, true);
    while (e) {
        if (e & 1) R = R * A;
        A = A * A; e >>= 1;
    }
    return R;
}
\end{lstlisting}

\section{数位 DP}

\cardmeta{现场常用}{统计 \([0,R]\) 中满足数字限制的数}{状态数乘位数}{前导零、tight 和 \(L-1\) 是三大边界}

\begin{lstlisting}[style=cpp]
string s;
// pos=当前位，tight=是否仍贴着上界，started=是否已经出现非前导零
int memo[20][2][2];
bool seen[20][2][2];

int dfs(int pos, int tight, int started) {
    // 到末尾代表构造出一个合法数字；这里把“0”也计入
    if (pos == (int)s.size()) return 1;
    if (seen[pos][tight][started]) return memo[pos][tight][started];
    seen[pos][tight][started] = true;
    int digit = s[pos] - '0';
    int lim = tight ? digit : 9, ans = 0;
    for (int d = 0; d <= lim; ++d) {
        if (d == 4) continue;
        // 只有仍贴着上界且 d 等于当前上界位时，下一层才 tight
        ans += dfs(pos + 1, tight && d == digit, started || d != 0);
    }
    return memo[pos][tight][started] = ans;
}
\end{lstlisting}

\begin{warnbox}
正式代码中 `tight \&\& d == lim` 应写成与当前上界数字比较，而不是与固定 9 比较。统计 \([L,R]\) 时使用 `F(R)-F(L-1)`，并单独处理 \(L=0\)。
\end{warnbox}

\section{SOS DP}

\cardmeta{现场常用}{子集和、超集和、位掩码统计}{\(O(n2^n)\)}{掩码位数通常不超过 20}

\begin{lstlisting}[style=cpp]
// f[mask] 初始为恰好 mask 的值；逐位把子集贡献累加到超集
for (int bit = 0; bit < n; ++bit)
    for (int mask = 0; mask < (1LL << n); ++mask)
        if (mask & (1LL << bit))
            f[mask] += f[mask ^ (1LL << bit)]; // 子集和
\end{lstlisting}

\section{Andrew 凸包}

\cardmeta{现场常用}{凸包、极值点、旋转卡壳前处理}{排序 \(O(n\log n)\)}{重复点先去重；保留边界点时改叉积条件}

\begin{lstlisting}[style=cpp]
struct Point {
    // 使用 long double 保存坐标，eps 来自 muban.cpp
    long double x, y;
    bool operator<(const Point &o) const {
        return x != o.x ? x < o.x : y < o.y;
    }
};
long double cross(Point a, Point b, Point c) {
    // (b-a) x (c-a)：>0 左转，<0 右转，约等于 0 共线
    return (b.x - a.x) * (c.y - a.y)
         - (b.y - a.y) * (c.x - a.x);
}
vector<Point> convex_hull(vector<Point> p) {
    // 先排序去重，再分别构造下凸壳和上凸壳
    sort(p.begin(), p.end());
    p.erase(unique(p.begin(), p.end(), [](Point a, Point b) {
        return fabsl(a.x - b.x) < eps && fabsl(a.y - b.y) < eps;
    }), p.end());
    if (p.size() <= 1) return p;
    vector<Point> h;
    for (auto x : p) {
        // 右转或共线时弹出中间点，保留外层边界
        while (h.size() >= 2 &&
               cross(h[h.size()-2], h.back(), x) <= eps) h.pop_back();
        h.push_back(x);
    }
    int lower = h.size();
    // lower 记录下凸壳结束位置，之后追加上凸壳
    for (int i = (int)p.size() - 2; i >= 0; --i) {
        while ((int)h.size() > lower &&
               cross(h[h.size()-2], h.back(), p[i]) <= eps) h.pop_back();
        h.push_back(p[i]);
    }
    h.pop_back();
    return h;
}
\end{lstlisting}
